\documentclass[11pt]{article}
\usepackage{amsmath,amssymb,amsthm,geometry,hyperref}
\geometry{margin=1in}
\allowdisplaybreaks
\setlength{\emergencystretch}{2em}
\newtheorem{theorem}{Theorem}
\newtheorem{proposition}{Proposition}
\newtheorem{lemma}{Lemma}
\title{Stolarsky monotonicity under logarithmic-mean smoothing\\
\large A solution of MathDB 376979}
\author{Research notes in Principia Math In Progress}
\date{11 September 2026}
\begin{document}
\maketitle
\begin{abstract}
We prove both inequalities in Conjecture 16 of Horwitz's \emph{Invariant Means}
for every positive triple, with strictness unless all entries coincide.
A fractional-linear normalization turns the question into two constrained
extremum problems. Boundary estimates and signed-integral identities reduce
both problems to triples with two equal entries. That remaining case follows
from an explicit positive-coefficient certificate. The live MathDB record
was checked on 11 September 2026 and reported open, with no posted claims.
\end{abstract}

\section{Task 1: the problem and its solution}
For $a,b>0$ put $L(a,b)=(b-a)/(\log b-\log a)$, with $L(a,a)=a$, and set
\[
 T(a,b,c)=(L(b,c),L(c,a),L(a,b)).
\]
Let $W=(W_1,W_2,W_3)$ be uniform on the two-simplex
$\{w_i\geq0:\sum w_i=1\}$, and define
\[
 F_q(a,b,c)=\mathbb E(W_1a+W_2b+W_3c)^{-q},\qquad
 M_q=F_q^{-1/q}\quad(q>0).
\]
The source's means are $U_0=M_2$ and $U_1=M_1$.
\begin{theorem}\label{main}
For every $a,b,c>0$,
\[
 M_2(T(a,b,c))\geq M_2(a,b,c),\qquad
 M_1(T(a,b,c))\leq M_1(a,b,c).
\]
Equality in either inequality holds if and only if $a=b=c$.
\end{theorem}
The theorem answers both assertions in the source's Conjecture 16, including
coincident coordinates. It also implies Conjecture 15's bounds for the
invariant logarithmic mean. We count the present record as one problem;
the latter consequence is not counted as a separate solve here.

\section{Integral formulation}
Writing $P(t)=(a+t)(b+t)(c+t)$, direct integration gives
\[
 F_2(a,b,c)=2\int_0^\infty\frac{dt}{P(t)},\qquad
 F_1(a,b,c)=2\int_0^\infty\frac{t\,dt}{P(t)}.
\]
These formulas remain nonsingular when coordinates coincide and are useful
for both analysis and reliable numerical evaluation.
For distinct coordinates they agree with the divided-difference formulas
\[
 F_2=-2[\log](a,b,c),\qquad F_1=2[x\log x](a,b,c).
\]
For completeness, integrating twice over the simplex gives
$\mathbb E(X+t)^{-3}=1/P(t)$, where $X=W_1a+W_2b+W_3c$.
Now use
$2\int_0^\infty(X+t)^{-3}\,dt=X^{-2}$ and
$2\int_0^\infty t(X+t)^{-3}\,dt=X^{-1}$ and Tonelli's theorem.

\section{A fractional-linear reduction of the general problem}
\begin{proposition}
Put $A=F_1(a,b,c)$, $B=F_2(a,b,c)$, $d=A/B$, and
$P(t)=(a+t)(b+t)(c+t)$. Then
\[
 T_i(a,b,c)=\frac{2}{A+B a_i}=\frac{2/B}{a_i+d}.
\]
The full pair of conjectured inequalities is equivalent to
\begin{align}
 P(d)\int_d^\infty\frac{dt}{P(t)}&\geq d,\label{tail-one}\\
 \left(\int_0^\infty\frac{dt}{P(t)}\right)P(d)
 \int_d^\infty\frac{t-d}{P(t)}\,dt&\leq1,\label{tail-two}
\end{align}
where $d=\int_0^\infty t/P(t)\,dt\big/\int_0^\infty1/P(t)\,dt$.
\end{proposition}
\begin{proof}
For $\{i,j,k\}=\{1,2,3\}$, the integral formulas yield
\[
 A+a_iB=2\int_0^\infty\frac{dt}{(t+a_j)(t+a_k)}
 =\frac{2}{L(a_j,a_k)}.
\]
The last identity also holds when $a_j=a_k$, by direct integration.
For a positive triple $r$, substitution $t=1/u$ shows
\[
 F_1(r^{-1})=(r_1r_2r_3)F_2(r),\qquad
 F_2(r^{-1})=(r_1r_2r_3)F_1(r).
\]
Homogeneity therefore gives
\[
 F_1(T)=\frac B2 P(d)F_2(a+d,b+d,c+d),\qquad
 F_2(T)=\frac{B^2}{4}P(d)F_1(a+d,b+d,c+d).
\]
Insert the integral formulas and $A=dB$ to obtain
\eqref{tail-one} from $F_1(T)\geq F_1(a,b,c)$ and
\eqref{tail-two} from $F_2(T)\leq F_2(a,b,c)$.
\end{proof}
The next sections prove the general inequalities through a different
normalization. The displayed tail inequalities are consequently true as well.

\section{Complete proof when two coordinates coincide}
\begin{proposition}\label{repeated}
If at least two of $a,b,c>0$ coincide, both conjectured inequalities hold.
Both are strict unless $a=b=c$.
\end{proposition}
\begin{proof}
By symmetry and homogeneity it suffices to consider $(1,1,x)$, $x>0$.
Write $s=\log x$ and $\ell=(e^s-1)/s$, with $\ell=1$ at $s=0$.
Then $T(1,1,x)$ is a permutation of $(1,\ell,\ell)$. Direct integration gives
\begin{align*}
 F_1(1,1,x)&=\frac{2(x\log x-x+1)}{(x-1)^2},&
 F_2(1,1,x)&=\frac{2(x-1-\log x)}{(x-1)^2},\\
 F_1(1,\ell,\ell)&=\frac{2(\ell-1-\log\ell)}{(\ell-1)^2},&
 F_2(1,\ell,\ell)&=\frac{2(\log\ell-1+1/\ell)}{(\ell-1)^2}.
\end{align*}
The two desired positive differences are identical:
\[
 F_1(1,\ell,\ell)-F_1(1,1,x)
 =F_2(1,1,x)-F_2(1,\ell,\ell)
 =\frac{2H(s)}{(\ell-1)^2},
\]
where
\[
 H(s)=\frac{(\ell-1)^3}{s\ell^2}-\log\ell+1-\frac1\ell.
\]
The apparent singularity at zero is removable and $H(0)=0$.
Differentiation gives
\[
 H'(s)=\frac{N(s,e^s)}{s^3(e^s-1)^3},
\]
with the following polynomial in $E$:
\begin{align*}
 N(s,E)={}&(s-2)E^4+(-s^3+s^2+8)E^3\\
 &+(s^4-2s^3-3s^2-6s-12)E^2\\
 &+(s^4+3s^3+3s^2+8s+8)E-s^2-3s-2.
\end{align*}
We now certify that every nonzero Taylor coefficient of each entire function
$N(s,e^s)$ and $-e^{4s}N(-s,e^{-s})$ is positive. It follows that
$N(s,e^s)>0$ for $s>0$, and $N(s,e^s)<0$ for $s<0$.
The denominator of $H'$ is positive on both half-lines. Thus $H$ decreases
to zero on the negative half-line and increases from zero on the positive
half-line. In particular $H(s)>0$ for $s\ne0$.

Here are complete algebraic certificates for the coefficient assertion.
Let $A_n$ and $B_n$ be the $n$th derivatives at zero of the respective
entire functions. They vanish for $0\leq n\leq8$, and direct integer
calculation gives
\[
\begin{array}{c|r|r}
n&A_n&B_n\\\hline
9&2520&2520\\
10&46200&54600\\
11&489720&674520\\
12&3945480&6255480\\
13&26865696&48487296\\
14&163315152&332155824\\
15&915554640&2079277200
\end{array}
\]
For $n\geq5$, writing $(n)_j=n(n-1)\cdots(n-j+1)$, expansion of the
four exponential terms gives
\begin{align*}
 A_n={}&4^{n-1}(n-8)
 +3^{n-3}\{-n(n-1)(n-5)+216\}\\
 &+2^{n-4}P_2(n)+(n)_4+3(n)_3+3(n)_2+8n+8,\\
 P_2(n)={}&n^4-10n^3+11n^2-50n-192.
\end{align*}
For $n\geq16$, $P_2(n)>0$, since
\[
P_2(m+16)=m^4+54m^3+1067m^2+9006m+26400.
\]
Also
\[
4^{n-1}(n-8)>3^{n-3}n(n-1)(n-5)
\]
at $n=16$, with difference $4380921872$.
The ratio of the left side to the right side increases
thereafter: the required numerator is
\[
 4(n-7)(n-1)(n-5)-3(n-8)(n+1)(n-4),
\]
whose coefficients in $m=n-16$ are $(1,29,288,1044)$.
This proves $A_n>0$ for every $n\geq16$.

For the second function the formula is
\[
 B_n=4^{n-2}P(n)-3^{n-4}Q(n)-2^{n-4}R(n)-S(n),
\]
where
\begin{align*}
P(n)&=n^2-13n+32,\\
Q(n)&=n^4-15n^3+65n^2-267n+648,\\
R(n)&=n^4-2n^3-13n^2+62n-192,\\
S(n)&=n^3-2n^2+n+8.
\end{align*}
All four polynomials have positive coefficients after $n=m+16$.
At $n=16$, the positive term exceeds the sum of the three subtracted
terms by $12156288144$. Each ratio of a subtracted term to the positive
term strictly decreases for $n\geq16$. To verify this last assertion,
the coefficient lists in $m=n-16$ of
\[
4P(n+1)Q(n)-3P(n)Q(n+1),\quad
4P(n+1)R(n)-2P(n)R(n+1),
\]
and $4P(n+1)S(n)-P(n)S(n+1)$, respectively, are
\begin{align*}
 &(1,64,1677,22884,169766,637568,912000),\\
 &(2,162,5410,95486,940476,4905664,10598400),\\
 &(3,200,5273,68572,438232,1094400).
\end{align*}
They are all positive. Hence $B_n>0$ for every $n\geq16$ as well.
This completes the coefficient certificate and the proof.
\end{proof}
The companion script \texttt{verify\_repeated\_case.py} regenerates the
derivative identity, the finite table, and all infinite-tail polynomial
certificates using exact rational arithmetic. This restricted theorem is the
input needed in the global extremum argument below; it is not counted separately.

\section{Normalization to a constraint in the unit cube}
For a positive triple $v$, abbreviate
\[
 A(v)=F_1(v),\qquad C(v)=F_2(v),\qquad
 P_v(t)=\prod_{i=1}^3(t+v_i).
\]
For $p\in(0,1)^3$, set $q=\boldsymbol1-p$ and
\[
 K=\prod_i p_i,\quad
 g(p)=C(p)-C(q),\quad D(p)=A(p)+A(q)+C(p),\quad
 J(p)=K C(p)D(p).
\]
Let $\mathcal S=\{p\in(0,1)^3:g(p)=0\}$.

\begin{lemma}\label{cube}
For $v_i=(1-p_i)/p_i$ one has
\[
 A(v)=K D(p),\qquad
 C(v)=K\{A(p)+A(q)+C(q)\}.
\]
Consequently Theorem~\ref{main} is equivalent to
\[
 A(p)\geq2,\qquad J(p)\leq4\quad(p\in\mathcal S).
\]
Both inequalities are strict except at
$p=(1/2,1/2,1/2)$.
\end{lemma}
\begin{proof}
For distinct coordinates, write $[f]_p$ for the second divided difference
of $f$ at the three coordinates of $p$. The change of nodes gives
\[
 [f]_v=K\,[x f((1-x)/x)]_p.
\]
Since $A=2[x\log x]$ and $C=-2[\log x]$, substituting these functions
and expanding $\log((1-x)/x)$ gives the two asserted identities.
They extend to all coordinates by continuity. In particular
$A(v)-C(v)=K g(p)$.

For an arbitrary original triple $a$, let $d=A(a)/C(a)>0$ and $v=a/d$.
Homogeneity gives $A(v)=C(v)=B$, say. Thus
$p_i=1/(1+v_i)$ belongs to $\mathcal S$ and $B=K D(p)$.
The fractional-linear formula proved earlier now reads
$T(v)=(2/B)p$. Hence
\[
 F_1(T(v))=\frac B2 A(p),\qquad
 F_2(T(v))=\frac{B^2}{4}C(p).
\]
The two desired inequalities are exactly $A(p)\geq2$ and $J(p)\leq4$.
Conversely every $p\in\mathcal S$ comes from the normalized triple
$v_i=(1-p_i)/p_i$. Equal coordinates correspond under this map, and the
only constant triple on $\mathcal S$ is $(1/2,1/2,1/2)$.
\end{proof}

\section{Boundary estimates}
\begin{lemma}\label{boundary}
Along every sequence in $\mathcal S$ approaching the boundary of the unit cube,
\[
 \liminf A(p)\geq2,\qquad \limsup J(p)\leq4.
\]
\end{lemma}
\begin{proof}
Write $a=\min_i p_i$ and $c=\max_i p_i$. If $a\to0$, then
\[
 C(p)\geq\frac12\int_0^1\frac{dt}{a+t}\longrightarrow\infty,
\]
because the other two factors of $P_p(t)$ are at most $2$ on $[0,1]$.
Since $C(q)=C(p)$, this forces $\min_i q_i=1-c\to0$; otherwise $C(q)$
would stay bounded. The converse follows by interchanging $p$ and $q$.
After passing to a convergent subsequence, every approach to the boundary
therefore has $a\to0$, $c\to1$, and $C(p)=C(q)\to\infty$.

Coordinate monotonicity gives $A(p)\geq A(a,1,1)\to2$, proving the first
bound. Also $A(p)/C(p)\to0$: for each fixed $\varepsilon>0$, eventually
$c\geq1/2$ and
\[
 \frac{A(p)}{C(p)}\leq\varepsilon+
 \frac{2}{C(p)}\int_\varepsilon^\infty\frac{dt}{t(t+1/2)}.
\]
The tail integral is finite. First take the limit and then let
$\varepsilon\downarrow0$. The same argument gives $A(q)/C(q)\to0$.
Thus $D(p)/C(p)\to1$.

If $a\leq b\leq c$ are the coordinates of $p$, then
\[
 C(p)\leq\frac2c\int_0^\infty\frac{dt}{(t+a)(t+b)}
 =\frac{2}{c L(a,b)}\leq\frac{2}{c\sqrt{ab}}.
\]
Here $L(a,b)\geq\sqrt{ab}$ follows immediately from
$\sinh u/u\geq1$. Consequently
\[
 K C(p)^2\leq\frac4c,
 \qquad J(p)=K C(p)^2\frac{D(p)}{C(p)}
 \quad\Longrightarrow\quad\limsup J(p)\leq4.
\]
The subsequence argument proves the assertions for arbitrary boundary
approaches, not just approaches to a specified face.
\end{proof}

\section{Eliminating extrema with three distinct coordinates}
\begin{lemma}\label{lower-extrema}
No constrained local minimum of $A$ on $\mathcal S$ has three distinct
coordinates.
\end{lemma}
\begin{proof}
Fix an interior point of $\mathcal S$, and write
$P(u)=\prod_i(u+p_i)$ and $Q(u)=\prod_i(u+q_i)$.
Differentiation under the integrals gives
\begin{align*}
 \partial_i A(p)&=-2\int_0^\infty\frac{t\,dt}{(t+p_i)P(t)},\\
 \partial_i g(p)&=-2\int_0^\infty
 \left\{\frac1{(t+p_i)P(t)}+\frac1{(t+q_i)Q(t)}\right\}dt<0.
\end{align*}
Thus $\nabla g\ne0$, and a constrained extremum satisfies
$\partial_i A=\lambda\partial_i g$, with $\lambda>0$ since both
derivatives are negative. Equivalently,
\[
 \int_0^\infty\frac{(t-\lambda)\,dt}{(t+p_i)P(t)}
 -\lambda\int_0^\infty\frac{dt}{(t+q_i)Q(t)}=0.
\]
In the second integral put $u=-1-t$. Its denominator becomes
$(u+p_i)P(u)$, since both factors change sign. With
$R_i(u)=P(u)/(u+p_i)$, the three equations become
\[
 \int R_i(u)\,d\mu(u)=0,
\]
where the signed measure, supported on $(-\infty,-1)\cup(0,\infty)$, has
density
\[
 \frac{d\mu}{du}=
 \begin{cases}
 -\lambda/P(u)^2,&u<-1,\\
 (u-\lambda)/P(u)^2,&u>0.
 \end{cases}
\]
If the $p_i$ are distinct, the three polynomials $R_i$ span all polynomials
of degree at most two: evaluation at $u=-p_j$ proves their linear independence.
In particular $\int(u-\lambda)\,d\mu(u)=0$. But its integrand is strictly
positive on $u<-1$ and nonnegative on $u>0$, with positive integral.
All these integrals converge absolutely. This contradiction proves the lemma.
\end{proof}

\begin{lemma}\label{upper-extrema}
No constrained local maximum of $J$ on $\mathcal S$ has three distinct
coordinates.
\end{lemma}
\begin{proof}
Use $f=\log J$ and the same $P,Q$ as above. The Lagrange equations give
$\nabla f=\lambda\nabla g$. We first show $\lambda=0$.
The homogeneity and uniform-shift derivative identities are
\[
 \sum_i v_i\partial_i A(v)=-A(v),\quad
 \sum_i v_i\partial_i C(v)=-2C(v),\quad
 \sum_i\partial_i A(v)=-C(v),\quad
 \sum_i\partial_i C(v)=-\frac2{\prod_i v_i}.
\]
The last two identities also follow directly by integration by parts.
Applying them to $q=\boldsymbol1-p$ gives, on $g=0$,
\[
 p\cdot\nabla D=-D,\qquad
 p\cdot\nabla f=3-2-1=0,\qquad
 p\cdot\nabla g=-\frac2{Q(0)}<0.
\]
Taking the scalar product of the Lagrange equations with $p$ therefore gives
$\lambda=0$, so all three partial derivatives of $f$ vanish.

Writing $C=C(p)$ and $D=D(p)$, these equations are
\[
 \frac1{p_i}
 -2\int_0^\infty
 \frac{1/C+(t+1)/D}{(t+p_i)P(t)}\,dt
 +\frac2D\int_0^\infty\frac{t}{(t+q_i)Q(t)}\,dt=0.
\]
The substitution $u=-1-t$ in the last integral now expresses them as
$\int R_i(u)\,d\nu(u)=0$, where
\[
 d\nu(u)=\frac1{P(0)}\,d\delta_0(u)
 +\begin{cases}
 \displaystyle\frac{2(-1-u)}{D P(u)^2}\,du,&u<-1,\\[4pt]
 \displaystyle-\frac{2\{1/C+(u+1)/D\}}{P(u)^2}\,du,&u>0,
 \end{cases}
\]
and $\delta_0$ denotes the unit point mass at zero. Indeed,
$R_i(0)/P(0)=1/p_i$. If the coordinates are distinct, the polynomials
$R_i$ span all polynomials of degree at most two. Thus $\int u\,d\nu(u)=0$.
In fact the integrand is strictly negative on both open intervals, while
the point mass contributes zero. The integral is absolutely convergent
and negative, a contradiction.
\end{proof}

\section{Completion of the global proof}
\begin{proof}[Proof of Theorem~\ref{main}]
Suppose $A(p)<2$ somewhere on $\mathcal S$. A minimizing sequence for $A$
has a subsequence converging in the closed cube. Lemma~\ref{boundary}
excludes convergence to the boundary when the infimum is less than $2$.
The infimum is therefore attained at an interior point of $\mathcal S$.
Lemma~\ref{lower-extrema} forces two coordinates to coincide, contradicting
Proposition~\ref{repeated} and the equivalence in Lemma~\ref{cube}.
Hence $A(p)\geq2$ on $\mathcal S$.

If $J(p)>4$ somewhere on $\mathcal S$, Lemma~\ref{boundary} similarly
forces its supremum to be attained in the interior. For clarity, it also
ensures this supremum is finite: unbounded values on a compactness sequence
could accumulate neither at the boundary nor at an interior point where
$J$ is continuous. Lemma~\ref{upper-extrema} then forces two equal
coordinates, again contradicting Proposition~\ref{repeated}.
Thus $J(p)\leq4$ throughout $\mathcal S$.

An interior equality point for either established bound is itself a
constrained extremum. The same lemmas force two equal coordinates there,
and the strictness in Proposition~\ref{repeated} forces all three to equal
$1/2$. Lemma~\ref{cube} and homogeneity transfer these conclusions to every
positive original triple, proving both inequalities and their equality cases.
\end{proof}

\paragraph{Consequence for the invariant mean.}
The source proves convergence of the logarithmic-mean iteration to
$(L_3,L_3,L_3)$. Along that iteration, the theorem makes $U_0$ increasing
and $U_1$ decreasing. Continuity therefore gives
$U_0(a,b,c)\leq L_3(a,b,c)\leq U_1(a,b,c)$.
If the original triple is nonconstant, the first step is strict, so both
bounds are strict as well.

\section{Local behavior for three distinct coordinates}
\begin{proposition}
For fixed $q>0$, put $a_i=1+u_i$, where $\sum_i u_i=0$. As $u\to0$,
\begin{align*}
 M_q(T(a))-M_q(a)
 ={}&\frac{3q-5}{96}\sum_i u_i^2
 +\frac{19-4q-3q^2}{480}\sum_i u_i^3
 +O(\lVert u\rVert^4).
\end{align*}
Consequently the two conjectured strict inequalities hold throughout
some neighborhood of the diagonal, except on the diagonal itself.
At $q=5/3$, the difference takes both signs arbitrarily close to the diagonal.
\end{proposition}
\begin{proof}
Uniform simplex moments satisfy
\[
 \mathbb E W_1^iW_2^jW_3^k=\frac{2i!j!k!}{(i+j+k+2)!}.
\]
In particular, for $Z=\sum_iW_i u_i$ we have
$\mathbb EZ=0$, $\mathbb EZ^2=\sum_i u_i^2/12$, and
$\mathbb EZ^3=\sum_i u_i^3/30$. Taylor expansion therefore gives
\[
 M_q(a)=1-\frac{q+1}{24}\sum_i u_i^2
 +\frac{(q+1)(q+2)}{180}\sum_i u_i^3+O(\lVert u\rVert^4).
\]
For the logarithmic means, direct expansion of their reciprocal integral is
\[
 L(1+v,1+w)=1+\frac{v+w}{2}-\frac{(v-w)^2}{12}
 +\frac{(v+w)(v-w)^2}{24}+O((|v|+|w|)^4).
\]
Write $\mu$ for the arithmetic mean of $T(a)$ and $v_i=T_i(a)-\mu$.
Using $\sum_i u_i=0$, the preceding expansion yields
\begin{align*}
 \mu&=1-\frac1{12}\sum_i u_i^2+\frac1{24}\sum_i u_i^3
 +O(\lVert u\rVert^4),\\
 \sum_i v_i^2&=\frac14\sum_i u_i^2-\frac14\sum_i u_i^3
 +O(\lVert u\rVert^4),\\
 \sum_i v_i^3&=-\frac18\sum_i u_i^3+O(\lVert u\rVert^4).
\end{align*}
Apply the expansion of $M_q$ at mean $\mu$, using homogeneity, and subtract
the expansion at $a$. This proves the asserted coefficients. Analyticity
on the positive octant makes the remainder uniform in direction near zero.
The quadratic coefficients at $q=1,2$ are nonzero with the required signs.
At $q=5/3$, the cubic term is $\sum_i u_i^3/120$; taking
$u=(t,t,-2t)$ and letting $t$ approach zero from either side proves
the final assertion.
\end{proof}
The exact symbolic coefficients are independently regenerated from the
defining integrals by \texttt{verify\_local\_expansion.py}.

\section{Follow-up tasks 2--5}
\begin{enumerate}
\renewcommand{\labelenumi}{\textbf{Task \arabic{enumi}.}}
\setcounter{enumi}{1}
\item Classify the exponents $q>0$ for which $M_q\circ T-M_q$ has a fixed
global sign. The local expansion proves that a positive sign requires
$q>5/3$ and a negative sign requires $q<5/3$; the critical exponent has
both signs locally. Determine the full admissible parameter ranges.
\item Determine the best constants in quantitative versions of the two
inequalities under $\max(a,b,c)/\min(a,b,c)\leq R$, normalizing each gap
by $\sum_i(a_i-\bar a)^2/\bar a$. Find their asymptotic behavior as
$R\downarrow1$ and $R\to\infty$.
\item Extend the theorem to $n$ variables. Define $M_q^{(n)}$ using the
uniform $(n-1)$-simplex and put
$(T_n(a))_i=M_1^{(n-1)}(a_1,\ldots,\widehat{a_i},\ldots,a_n)$.
Determine which exponents give global monotonicity under $T_n$, and
whether the extremum reduction still reduces to repeated coordinates.
\item Derive effective, certified error bounds for iteration of $T$ toward
the invariant logarithmic mean. Determine the sharp asymptotic contraction
rate of the interval with endpoints $U_0(T^k(a,b,c))$ and
$U_1(T^k(a,b,c))$, and give a stopping rule for any prescribed accuracy.
\end{enumerate}

\begin{thebibliography}{9}
\bibitem{horwitz} A. Horwitz, \emph{Invariant Means}, arXiv:math/0007095v1
(2000), Conjecture 16; journal version, J. Math. Anal. Appl. 270 (2002),
499--518. \url{https://arxiv.org/abs/math/0007095}.
\bibitem{mathdb} MathDB, problem 376979, live record retrieved 11 September
2026. \url{https://mathdb.com/p/376979}.
\end{thebibliography}
\end{document}
